\subtitle{\bf Aula 21 - Introdução}
\begin{frame}[t,plain]
    \titlepage
\end{frame}

\begin{frame}{Objetivos}

    \begin{itemize}
        \item
        Distinguir verificação de tipos de inferência de tipos e apresentar a
        linguagem MiniML.
        \item
        Apresentar a abordagem de geração e resolução de restrições para
        inferência de tipos, incluindo o algoritmo de unificação de Robinson e
        a generalização de Milner.
        \item
        Apresentar a eliminação de tipos como ponte entre o inferidor e o
        backend do compilador, e o pipeline completo MiniML -> Lambda -> TImp.
    \end{itemize}
\end{frame}


\section{Introdução}\label{introduuxe7uxe3o}}

\begin{frame}{Verificação vs.~Inferência}

    \textbf{Verificação de tipos}: dado um programa \textbf{com anotações},
    decide se é bem tipado.

    \textbf{Inferência de tipos}: dado um programa \textbf{sem anotações},
    encontra os tipos mais gerais --- ou reporta que nenhum existe.

    Regra (T-Abs) no STLC:

    \[\frac{\Gamma,\, x : T_1 \vdash t : T_2}{\Gamma \vdash \lambda x {:} T_1.\; t : T_1 \to T_2}\]

    \begin{itemize}

        \item
        Verificação: \(T_1\) é fornecido pelo programador
        \item
        Inferência: \(T_1\) deve ser \textbf{descoberto} a partir do uso de
        \(x\) no corpo
    \end{itemize}
\end{frame}

\begin{frame}[fragile]{A Linguagem MiniML}

    \begin{lstlisting}
        e ::= x              -- variável
        | ℓ              -- literal (int ou bool)
        | e₁ e₂          -- aplicação
        | λx. e          -- abstração sem anotação
        | let x = e₁ in e₂  -- definição polimórfica
    \end{lstlisting}

    \begin{itemize}

        \item
        Abstrações \textbf{não} têm anotações de tipo
        \item
        O inferidor deduz tudo automaticamente
    \end{itemize}
\end{frame}

\begin{frame}{Tipos e Esquemas}

    \textbf{Tipos monomórficos:}

    \[T ::= \mathbf{Int} \mid \mathbf{Bool} \mid T \to T \mid \alpha\]

    \textbf{Esquemas de tipo} (tipos polimórficos):

    \[\sigma ::= T \mid \forall \vec{\alpha}.\; T\]

    \begin{itemize}

        \item
        Esquema \(\forall \alpha_1 \ldots \alpha_n.\; T\) generaliza \(T\) nas
        variáveis \(\alpha_i\)
        \item
        A \textbf{instanciação} substitui variáveis quantificadas por tipos
        novos
        \item
        A \textbf{generalização} de \(T\) em \(\Gamma\):
        \(\forall \vec{\alpha}.\; T\) onde
        \(\vec{\alpha} = \mathrm{fv}(T) \setminus \mathrm{fv}(\Gamma)\)
    \end{itemize}
\end{frame}


\section{Linguagem de Restrições}\label{linguagem-de-restriuxe7uxf5es}}

\begin{frame}{Gramática de Restrições}

    \[\begin{array}{rcll}
    C & ::= & \top & \text{trivial} \\
    & \mid & T_1 = T_2 & \text{equação de tipos} \\
    & \mid & C_1 \wedge C_2 & \text{conjunção} \\
    & \mid & \exists \alpha.\; C & \text{variável existencial} \\
    & \mid & \mathbf{def}\; x : T\; \mathbf{in}\; C & \text{variável monomórfica} \\
    & \mid & \mathbf{let}\; x : T\; [C_1]\; \mathbf{in}\; C_2 & \text{definição polimórfica} \\
    & \mid & \mathbf{inst}(x,\, T) & \text{instanciação}
\end{array}\]

\begin{itemize}

    \item
    \(\exists \alpha.\; C\) --- variável de tipo local, visível apenas em
    \(C\)
    \item
    \(\mathbf{def}\) --- parâmetro de lambda (não polimórfico)
    \item
    \(\mathbf{let}\) --- polimorfismo de let (generaliza após resolver
    \(C_1\))
\end{itemize}
\end{frame}

\begin{frame}{Geração de Restrições \(\mathcal{G}(e, \alpha)\)}

    \(\mathcal{G}(e, \alpha)\) produz uma restrição cuja solução torna
    \(\alpha\) o tipo de \(e\):

    \[\mathcal{G}(x,\; \alpha) = \mathbf{inst}(x,\, \alpha) \quad\text{(G-Var)}\]

    \[\mathcal{G}(\ell,\; \alpha) = \alpha = \mathrm{tylit}(\ell) \quad\text{(G-Lit)}\]

    \[\mathcal{G}(e_1\; e_2,\; \alpha) = \exists \beta.\;\bigl(\mathcal{G}(e_1,\; \beta \to \alpha) \wedge \mathcal{G}(e_2,\; \beta)\bigr) \quad\text{(G-App)}\]

    \[\mathcal{G}(\lambda x.\; e,\; \alpha) = \exists \beta_1.\; \exists \beta_2.\; \bigl(\mathbf{def}\; x : \beta_1\; \mathbf{in}\; \mathcal{G}(e,\, \beta_2)\bigr) \wedge (\alpha = \beta_1 \to \beta_2) \quad\text{(G-Lam)}\]

    \[\mathcal{G}(\mathbf{let}\; x = e_1\; \mathbf{in}\; e_2,\; \alpha) = \exists \beta.\; \mathbf{let}\; x : \beta\; [\mathcal{G}(e_1,\, \beta)]\; \mathbf{in}\; \mathcal{G}(e_2,\, \alpha) \quad\text{(G-Let)}\]
\end{frame}

\begin{frame}{Exemplo de Geração}

    Para \(\lambda f.\; f\; \mathbf{true}\) com alvo \(\alpha_0\):

    \[\mathcal{G}(\lambda f.\; f\; \mathbf{true},\; \alpha_0) = \exists \beta_1\, \beta_2\, \beta_3.\; \bigl(\mathbf{def}\; f : \beta_1\; \mathbf{in}\; \mathbf{inst}(f,\, \beta_3 \to \beta_2) \wedge \beta_3 = \mathbf{Bool}\bigr) \wedge (\alpha_0 = \beta_1 \to \beta_2)\]

    A resolução força \(\beta_1 = \mathbf{Bool} \to \beta_2\).

    Tipo principal: \(\alpha_0 = (\mathbf{Bool} \to \beta_2) \to \beta_2\),
    generalizado em \(\beta_2\):

    \[\forall \beta.\; (\mathbf{Bool} \to \beta) \to \beta\]
\end{frame}


\section{O Resolvedor de
Restrições}\label{o-resolvedor-de-restriuxe7uxf5es}}

\begin{frame}{Unificação de Robinson}

    \(\mathrm{mgu}(T_1, T_2)\) encontra o \textbf{unificador mais geral}:

    \[\begin{array}{rcll}
    \mathrm{mgu}(\alpha, \alpha) & = & \varepsilon & \text{variável consigo mesma} \\
    \mathrm{mgu}(\alpha, T) & = & [\alpha \mapsto T] & \text{se } \alpha \notin \mathrm{fv}(T) \\
    \mathrm{mgu}(T_1 \to T_2,\; T_1' \to T_2') & = & S_2 \circ S_1 & S_1 = \mathrm{mgu}(T_1, T_1'),\; S_2 = \mathrm{mgu}(S_1(T_2), S_1(T_2')) \\
    \mathrm{mgu}(T, T) & = & \varepsilon & \text{tipos concretos iguais} \\
    \mathrm{mgu}(T_1, T_2) & = & \text{erro} & \text{caso contrário}
\end{array}\]

\textbf{Occurs check}: \(\alpha \notin \mathrm{fv}(T)\) --- evita tipos
infinitos como \(\alpha = \alpha \to \alpha\).
\end{frame}

\begin{frame}[fragile]{O Algoritmo de Resolução}

    \begin{lstlisting}[language=Haskell]
        solve :: Constraint -> Solve ()
        solve CTrue = pure ()
        solve (t1 :=: t2) = do
        s  <- askSubst
        s' <- mgu (apply s t1) (apply s t2)
        extSubst s'
        solve (c1 :&: c2) = solve c1 >> solve c2
        solve (CExists c) = do
        v <- freshTyVar
        solve (c v)
        solve (CDef x t c) =
        withLocalCtx x (Forall [] t) (solve c)
        solve (CLet x t c1 c2) = do
        solve c1
        scheme <- generalize t
        withLocalCtx x scheme (solve c2)
        solve (CInst x t) = do
        t' <- lookupVar x
        solve (t :=: t')
    \end{lstlisting}
\end{frame}


\section{Correção e Completude}\label{correuxe7uxe3o-e-completude}}

\begin{frame}{Teoremas Fundamentais}

    \textbf{Teorema (Correção)}: Se o resolvedor aceita \(C_e\) com
    substituição \(S\), então \(\vdash e : S(\alpha_0)\).

    \emph{O tipo produzido é válido segundo Hindley-Milner.}

    \textbf{Teorema (Completude)}: Se \(\Gamma \vdash e : T\), então o
    resolvedor aceita \(C_e\) e produz um tipo do qual \(T\) é instância.

    \emph{Se \(e\) é tipável, o algoritmo encontra o tipo mais geral.}

    \textbf{Teorema (Tipo Principal)}: Existe \(\sigma\) tal que: 1.
    \(\vdash e : T\) para todo \(T\) instância de \(\sigma\) 2. Todo tipo
    válido \(T'\) é instância de \(\sigma\)

    O algoritmo sempre encontra o \textbf{tipo principal}.
\end{frame}


\section{Eliminação de Tipos e
Pipeline}\label{eliminauxe7uxe3o-de-tipos-e-pipeline}}

\begin{frame}{Eliminação de Tipos (Type Erasure)}

    Tipos são necessários para análise estática, mas \textbf{desnecessários
    em execução}.

    A função \(\lfloor \cdot \rfloor : \mathit{TyExp} \to \mathit{Term}\)
    descarta anotações:

    \[\begin{array}{rcl}
    \lfloor x : T \rfloor & = & x \\
    \lfloor (te_1\; te_2) : T \rfloor & = & \lfloor te_1 \rfloor\; \lfloor te_2 \rfloor \\
    \lfloor \lambda x {:} T_1.\; te : T \rfloor & = & \lambda x.\; \lfloor te \rfloor \\
    \lfloor \mathbf{let}\; x = te_1\; \mathbf{in}\; te_2 : T \rfloor & = & (\lambda x.\; \lfloor te_2 \rfloor)\; \lfloor te_1 \rfloor
\end{array}\]

\textbf{let} -> beta-redex: \((\lambda x.\; t_2)\; t_1\) ---
semanticamente equivalente.
\end{frame}

\begin{frame}[fragile]{O Pipeline Completo de MiniML}

    \begin{lstlisting}[language=Haskell]
        processTImp :: String -> IO ()
        processTImp input =
        case parser input of
        Left err  -> putStrLn $ "Parse error: " ++ err
        Right ast ->
        case inferElab ast of
        Left err -> putStrLn $ "Type error: " ++ err
        Right (_, _, te, ty) -> do
        putStrLn $ "val : " ++ pretty ty    -- tipo exibido
        let lterm = erase te                -- elimina tipos
        prog  = lambdaToTImp lterm      -- closure conversion
        result <- interpret prog
        ...
    \end{lstlisting}

    \textbf{Quatro fases}: Parser -> Inferência/Elaboração -> Eliminação ->
    TImp
\end{frame}

\begin{frame}[fragile]{Exemplo Completo}

    Para \texttt{let f = \\x . x in f 42}:

    \begin{enumerate}
        \item
        \textbf{Inferência}: tipo \texttt{Int}, árvore tipada
        \(\mathbf{let}\; f = (\lambda x{:}\alpha.\; x){:}\alpha \to \alpha\; \mathbf{in}\; (f\; 42){:}\mathbf{Int}\)
        \item
        \textbf{Eliminação}: \((\lambda f.\; f\; 42)\; (\lambda x.\; x)\)
        \item
        \textbf{Closure conversion}: duas funções promovidas,
        \texttt{lam\_0} (identidade) e
        \texttt{lam\_1} (aplicação de
        \texttt{f})
        \item
        \textbf{Execução TImp}: imprime \texttt{42}
    \end{enumerate}
\end{frame}


\section{Conclusão}\label{conclusuxe3o}}

\begin{frame}{Sumário do Capítulo}

    \begin{itemize}

        \item
        \textbf{Inferência de tipos} resolve um problema mais difícil que
        verificação: sem anotações
        \item
        \textbf{Geração de restrições} \(\mathcal{G}(e, \alpha)\) transforma
        inferência em sistema de equações
        \item
        \textbf{Unificação de Robinson} encontra o unificador mais geral
        \item
        \textbf{Polimorfismo de let}: generalização após resolver restrições
        do corpo
        \item
        \textbf{Eliminação de tipos}: descarta anotações, reutiliza pipeline
        do \(\lambda\)-cálculo
    \end{itemize}
\end{frame}

\begin{frame}{Próximos Passos}

    \begin{itemize}

        \item
        \textbf{Subtipagem} (capítulo 19): \(S <: T\) --- todo \(S\) pode ser
        usado onde \(T\) é esperado
        \item
        \textbf{Tipos dependentes}: tipos que dependem de valores
        \item
        \textbf{Sistemas de efeitos}: rastrear efeitos colaterais no sistema
        de tipos
    \end{itemize}
\end{frame}
